\chapter{Introduction}%
\label{chp:introduction}
The application of Artificial Intelligence (AI) and computer vision in the field of sports science has grown significantly, offering new avenues for performance analysis beyond traditional, subjective coaching methods. This Final Year Project (FYP) introduces a proof-of-concept system designed to leverage these technologies for the biomechanical analysis of competitive swimming. The system aims to provide objective, data-driven insights by detecting key stroke phases and quantifying limb symmetry in freestyle swimming footage.

The primary goals of this project are twofold: firstly, to develop a robust method for segmenting freestyle swimming motion into biomechanically relevant phases, such as the catch, pull, push, and recovery. Secondly, the project aims to establish a methodology for analyzing and quantifying the bilateral symmetry of a swimmer's limb movements. The approach used will involve applying an off-the-shelf pose estimation model to extract keypoint data from video footage. This keypoint sequence will then be processed using machine learning techniques to achieve the aforementioned goals.

A core assumption of this project is that existing datasets of swimming footage, such as those from the DIVE and SWIM-360 projects, can be used for initial development and testing. Furthermore, as a proof-of-concept, the system's final evaluation will be primarily qualitative, relying on annotated samples and feedback from experienced swimming coaches.

At a high level, the system will intake raw swimming video footage. It will then apply a pre-trained pose estimation model to identify and track the coordinates of a swimmer’s body joints. These keypoint data will form the basis for two subsequent analyses: a sequence modeling approach for stroke phase detection, and a comparative analysis of left and right limb trajectories to quantify symmetry. The final output will be a visual overlay on the video, providing immediate feedback on stroke phases and symmetry metrics.

This report is structured as follows: Chapter 1 introduces the project's context, aims, and methodology. Chapter 2 provides the necessary background information and a review of relevant literature. Chapter 3 details the system's design and architecture. Chapter 4 explains the implementation process. Chapter 5 presents the evaluation and testing of the prototype. Finally, Chapter 6 discusses future work and concludes the report.